\section{The packet as a weighted ideal sum}\label{sec:packets}

We first identify the packet measure exactly, including its signs and
orbit volumes. Throughout the proof, $K$ is totally real of the fixed
prime degree $n$, and
\[
 R=\cO_K,\quad E=R^\times,\quad D=|\Disc(K)|,\quad
 \kappaK=\operatorname*{Res}_{s=1}\zeta_K(s).
\]
The absolute norm $\N I$ of a nonzero fractional $R$-ideal is extended
multiplicatively from integral ideals. Thus
$\operatorname{covol}(\sigma(I))=\sqrt D\,\N I$.
We use standard ideal theory and the unit theorem
\cite[Chapters~3--5]{Milne2020}, together with the analytic
class-number formula \cite[Theorem~19.12]{Sutherland2021}.
The classical packet interpretation is discussed in
\cite[Sections~2 and~5]{ELMV2011}.

\subsection{Saturation and local choices}

Let $J_0=RM$, $q=[J_0:M]$, and $Q=q\sqrt D$. For a rational prime
$\ell$, write $(J_0)_\ell=j_\ell R_\ell$, and put
$L_\ell=j_\ell^{-1}M_\ell$. Then
\begin{equation}\label{eq:local-normalization}
 L_\ell\subset R_\ell,\qquad R_\ell L_\ell=R_\ell,\qquad
 [R_\ell:L_\ell]=q_\ell,\qquad q=\prod_\ell q_\ell.
\end{equation}
Set $S=\{\ell:\ell\mid q\}$; off $S$, $L_\ell=R_\ell$.
Since $qJ_0\subset M$, one has $qR\subset\cO(M)$.
Conversely, a multiplier of $M$ has an integral endomorphism matrix
on $M$ and hence is an algebraic integer. It follows that
\begin{equation}\label{eq:discriminant-scale}
 qR\subset\cO(M)\subset R,\qquad
 D(M)\le Dq^{2n}\le Q^{2n}.
\end{equation}
Only the last implication, $D(M)\to\infty\Rightarrow Q\to\infty$,
will be needed.

These quantities are constant within the prescribed local type.
Indeed, if $M'_\ell=c_\ell M_\ell$, multiplication by $c_\ell$
identifies $(RM)_\ell/M_\ell$ with $(RM')_\ell/M'_\ell$.
It also identifies the local multiplier orders, since multiplication
in $K\otimes_\Q\Q_\ell$ is commutative. The multiplier order commutes
with completion: in a $\Z$-basis of $M$, its defining condition is
that the matrix of multiplication have integral entries, and the
same condition over $\Z_\ell$ defines the local order. Thus $q$,
$Q$, $\cO(M)$, and $D(M)$ are unchanged by local homothety.
Different local types can nevertheless have the same multiplier order.

Choose one fractional ideal $I$ in each ordinary ideal class of $R$.
Let $\mathcal N(I)$ be the set of lattices $N$ locally homothetic to $M$
with $RN=I$. If $I_\ell=i_\ell R_\ell$, its possible completions are
exactly
\begin{equation}\label{eq:local-choices}
 N_\ell=i_\ell z_\ell L_\ell,\qquad z_\ell\in R_\ell^\times.
\end{equation}
Indeed, taking maximal-order spans in a local homothety forces the
remaining multiplier to be a unit. Conversely, these local lattices
lie between $qI_\ell$ and $I_\ell$ and patch uniquely by the primary
decomposition of $I/qI$. Thus $\mathcal N(I)$ is finite,
$[I:N]=q$ for all its members, and its cardinality
$T=|\mathcal N(I)|$ is independent of $I$. The label count $T$ need
not equal the saturation index $q$. Local unit Haar probability induces
uniform probability on each finite set of local choices, so the choices at different
primes are independent. Every global homothety class in the
packet occurs in these sets, by moving its maximal-order span
to the chosen ideal-class representative.

\subsection{Real parameters and measure}

Let
\[
 C=\{t\in(\R^\times)^n:\textstyle\prod_j|t_j|=1\}.
\]
Give $C$ counting measure on its $2^n$ sign components and measure
$du_1\cdots du_{n-1}$ on the coordinates $u_j=\log|t_j|$.
The embedding of $E$ is a discrete cocompact subgroup of $C$.
If $F$ is a measurable fundamental domain, the unit theorem and
the totally real case of the analytic class-number formula
\cite[Theorem~19.12]{Sutherland2021} give
\begin{equation}\label{eq:unit-volume}
 \vol(F)=2^{n-1}R_K,\qquad
 h_K\vol(F)=\sqrt D\,\kappaK,
\end{equation}
where $h_K$ is the ordinary class number and $R_K$ the regulator
in deleted logarithmic coordinates. The factor $2^{n-1}$ comes
from the $2^n$ signs and the logarithmic kernel $\{1,-1\}$.
In the class-number formula this is the factor $2^n/w_K$, since
the roots of unity in a totally real field are exactly $\{1,-1\}$.

\begin{lemma}\label{lem:packet-model}
The maps
\[
 (N,t)\longmapsto (Q\N I)^{-1/n}\sigma(N)t
 \qquad (N\in\mathcal N(I),\ t\in C)
\]
induce a bijection from the disjoint union of the quotients by
\[
 \epsilon:(N,t)\longmapsto
       (\epsilon N,\sigma(\epsilon)^{-1}t),\qquad \epsilon\in E,
\]
onto the union of the packet orbits.
Counting measure on the $N$ labels times the stated measure on $C$
descends to the sum of the orbit Haar measures. In particular,
the packet probability is obtained by averaging uniformly over $I$
and $\mathcal N(I)$ and integrating over $F$ with measure
$dt/\vol(F)$.
\end{lemma}

\begin{proof}
Surjectivity follows from \eqref{eq:local-choices}. A global
homothety changes a normalized embedded lattice by an element
of $C$. Since $n$ is odd, each pair of opposite sign patterns has
exactly one representative in $W_n$; a common negative sign
does not change a lattice.

For injectivity, suppose two displayed normalized lattices coincide.
Absorb the ratio of the normalizing scalars into the real diagonal
map carrying $\sigma(N)$ onto $\sigma(N')$. Applying this map
to one nonzero $\alpha\in N$ gives $\sigma(\beta)$ for some
$\beta\in N'$. None of the coordinates of $\sigma(\alpha)$
vanishes, so the map is multiplication by $\sigma(\beta/\alpha)$.
Thus $N'=bN$ for $b\in K^\times$. Their maximal-order spans
represent the same ideal class and therefore equal the same
chosen $I$. Hence $bI=I$, so $b\in E$, and the parameters differ
by exactly the stated action.

Combine an $N$ label with a sign component of $C$. Its stabilizer in $E$
is the totally positive subgroup of
$\{\epsilon\in E:\epsilon N=N\}$, and has finite index in $E$.
It acts faithfully by logarithmic translations. By injectivity,
these are all periods of the corresponding $A_n$-orbit.
Each orbit of labels therefore contributes exactly one compact
$A_n$-orbit, with precisely its prescribed Haar volume.
This proves both finiteness and the measure assertion.
One may equally use $\mathcal N(I)\times F$ as a fundamental
domain, since the $E$-action preserves counting measure times $dt$.
The total measure of these fundamental domains is
$h_KT\vol(F)$. Dividing by this quantity gives precisely the
uniform averages over ideal classes and finite labels, followed
by integration against $dt/\vol(F)$.
\end{proof}

\subsection{Containment probabilities and unfolding}

Write $\mathfrak p\mid\ell$ for the primes of $R$ over $\ell$,
with ramification indices $e_{\mathfrak p}$ and residue degrees
$f_{\mathfrak p}$. Normalize $v_{\mathfrak p}$ by a uniformizer,
so $v_{\mathfrak p}(\ell)=e_{\mathfrak p}$ and
$\sum_{\mathfrak p\mid\ell}e_{\mathfrak p}f_{\mathfrak p}=n$.
For a tuple $\mathbf k_\ell=(k_{\mathfrak p})_{\mathfrak p\mid\ell}$
of nonnegative integers, define
\[
 P_\ell(\mathbf k_\ell)
 =\Pr_{z\in R_\ell^\times}\{x\in zL_\ell\},
\]
where $x$ has those component valuations and the probability is
normalized unit Haar measure. The definition is independent of $x$
because the units act transitively on each valuation shell.
For an integral nonzero $R$-ideal $\mathfrak a$, set
\[
 P(\mathfrak a)=\prod_{\ell\in S}
       P_\ell((v_{\mathfrak p}(\mathfrak a))_{\mathfrak p\mid\ell}).
\]
For $\alpha\in I\setminus\{0\}$, the probability that a uniform
$N\in\mathcal N(I)$ contains $\alpha$ is
$P((\alpha)I^{-1})$.

For a nonnegative Borel function $f$ on $\R^n$, put
$E_f(\Lambda)=\sum_{0\ne v\in\Lambda}f(v)$.
The product-fiber integral below integrates $f$ over the $C$-orbit
with absolute coordinate product $z$, using logarithmic Haar measure
and summing over all sign components:
\begin{equation}\label{eq:product-fiber}
 \cV_f(z)=\sum_{\omega\in\{\pm1\}^n}
 \int_{\R^{n-1}}
 f\left(\omega_1e^{u_1},\ldots,\omega_{n-1}e^{u_{n-1}},
      \omega_n z e^{-\sum_{j<n}u_j}\right)
 \dd u_1\cdots\dd u_{n-1},\qquad z>0.
\end{equation}

\begin{proposition}\label{prop:unfolding}
Writing $\mu_*=\mu_{K,M,\sigma}$, one has
\begin{equation}\label{eq:unfolding}
 \int_{X_n}E_f\dd\mu_*
 =\frac1{\sqrt D\,\kappaK}
       \sum_{\mathfrak a\subset R,\ \mathfrak a\ne0}
           P(\mathfrak a)\cV_f(\N\mathfrak a/Q).
\end{equation}
The identity holds for nonnegative extended integrals.
\end{proposition}

\begin{proof}
Lemma~\ref{lem:packet-model} expresses the left side as
\[
 \frac1{h_K\vol(F)}
 \sum_I\int_F\sum_{\alpha\in I\setminus\{0\}}
   P((\alpha)I^{-1})
   f((Q\N I)^{-1/n}\sigma(\alpha)t)\dd t.
\]
The map $\alpha\mapsto(\alpha)I^{-1}$ runs over the integral
ideals in the inverse class of $I$, each fiber being one free
$E$-orbit. Summing that orbit and integrating over $F$ unfolds
to an integral over $C$, by Tonelli. The absolute coordinate
product of the normalized vector is
$\N((\alpha)I^{-1})/Q$. Translating the first $n-1$ logarithms
and relabeling the signs identifies the $C$-integral with
\eqref{eq:product-fiber}, with Jacobian one.
Every integral ideal occurs in exactly one inverse class.
Equation~\eqref{eq:unit-volume} supplies the coefficient in
\eqref{eq:unfolding}.
\end{proof}

When $M$ is an $R$-ideal, $q=1$ and $P(\mathfrak a)=1$.
Thus \eqref{eq:unfolding} recovers the maximal-order formula.
